% Standalone TRACE four-layer architecture figure for the ITLT variant.
% Build: xelatex arch_trace.tex  → arch_trace.pdf  (then convert to PNG if needed).
\documentclass[border=6pt]{standalone}
\usepackage{fontspec}
\usepackage{polyglossia}
\setmainlanguage{ukrainian}
\setotherlanguage{english}
\setmainfont{DejaVu Serif}
\setsansfont{DejaVu Sans}
\setmonofont{DejaVu Sans Mono}
\usepackage{amsmath,amssymb,array}
\usepackage{tikz}
\usetikzlibrary{arrows.meta,positioning}
\begin{document}
\begin{tikzpicture}[
  font=\small,
  box/.style={draw, rounded corners, align=center, inner sep=5pt, minimum width=14cm}
]
  \node[box, fill=blue!8] (l4) {\textbf{L4 --- Обмежений нагляд людини}\\\footnotesize перевизначення будь-якого автоматичного рішення адміністратором \;|\; журнал дій};
  \node[box, below=4mm of l4, fill=orange!12] (l3) {\textbf{L3 --- Політика ескалації}\\\footnotesize пороги $\theta_{\text{зл}}$ і $\theta_{\text{пер}}$ \;|\; негайна відмова від злиття за конфлікту DOI \;|\; умови винесення на ручну перевірку};
  \node[box, below=4mm of l3, fill=green!8] (l2) {\textbf{L2 --- Навчені компоненти}\\[3pt]
    \footnotesize
    \begin{tabular}{@{}>{\raggedright\arraybackslash}p{0.46\linewidth}@{\hspace{5mm}}>{\raggedright\arraybackslash}p{0.46\linewidth}@{}}
      \textbf{L2a --- класичне машинне навчання}\newline оцінка зіставлення: близькість назв, відстань Левенштейна, коефіцієнт Жаккара; логістична функція $\sigma(\boldsymbol{\beta}\cdot\boldsymbol{\phi})$ &
      \textbf{L2b --- мовна модель як валідатор}\newline тлумачення природномовних запитів, пояснення, перевірка узгодженості \\
    \end{tabular}};
  \node[box, below=4mm of l2, fill=gray!12] (l1) {\textbf{L1 --- Детерміноване ядро} \,(якір довіри, без генерування)\\\footnotesize канонізація DOI \;|\; жорсткі обмеження (охоронець конфлікту DOI) \;|\; фіксація походження даних (W3C~PROV)};

  \draw[-{Stealth[length=2.6mm]}, thick] (l1.north) -- (l2.south);
  \draw[-{Stealth[length=2.6mm]}, thick] (l2.north) -- (l3.south);
  \draw[-{Stealth[length=2.6mm]}, thick] (l3.north) -- (l4.south);
\end{tikzpicture}
\end{document}
